\section{Introduction}

% Fala de iot, seus problemas e requisitos de segurança, explica por que é uma boa fazer analise de trafego
The Internet of Things (IoT) refers to the ubiquitous interconnection of smart devices and sensors in interoperable networks \cite{lin:2017}. Despite the resource limitations of these devices, IoT integration with edge/cloud computing paradigms is expanding its presence to a plethora of daily applications. This increases the volume of information exchange and, consequently, the attack surface \cite{tsimenidis:2022}. As a result, network traffic analysis plays a fundamental role in IoT cybersecurity \cite{hussain:2020}. By analyzing network traffic, it is possible to identify abnormal behaviors in the network by monitoring communication patterns and data exchanges. Network anomaly detection (NAD) proactively prevents malicious activities, such as cyberattacks, intrusions, or policy violations \cite{wang:2025, miguel:2026}. As IoT networks continue to grow in size and complexity, automated and intelligent approaches for traffic analysis that comply with their computational limitation have become necessary \cite{tsimenidis:2022}.

To analyze network flows efficiently, several techniques based on artificial intelligence (AI) and machine learning have been proposed \cite{folino:2021, alghushairy:2024, talukder:2024}. Among them, the Deep Neural Network (DNN) is currently the top choice for NAD, given its ability to model nonlinear pattern boundaries, better complying with the characteristics of network traffic. Some of its instances can spatially locate or address the temporal context of the anomaly, making them suitable for real-time Intrusion Detection Systems \cite{anis:2025}. 

Although DNNs achieve high accuracy results, they are also very resource-intensive in memory and computation. The high number of multiply-accumulate operations used for weight matrix calculations requires not only memory, but a huge amount of energy \cite{susskind:2022}. This proves impractical in resource-constrained and latency-sensitive networks \cite{hussain:2020, lello:2020, kalpani:2026}, such as IoT networks or Local Area Networks (LANs). Addressing this issue, weightless neural networks (WNN), such as the WiSARD (Wilkie, Stonham \& Aleksander’s Recognition Device) model, are being revisited as a promising alternative for low-power requirements \cite{susskind:2022}. WNN is based on Random Access Memory (RAM) nodes rather than traditional weight matrices, with entries coded as Boolean functions in lookup tables \cite{aleksander:2009, lello:2020, susskind:2022}. This approach enables lightweight learning, making it well-suited to the constraints of IoT networks.

Despite the benefits of WNN, most of its published applications are limited to image detection \cite{di:2020}, with scarce existing literature exploring cybersecurity domains. This is mostly due to the classification performance and dataset size requirements when compared with DNN \cite{miranda:2022, susskind:2022}. Also, it raises challenges regarding the encoding strategies. Specifically in NAD, these challenges concern the encoding of flow data into binary representations suitable for this type of model \cite{kappaun:2016}.

To address these challenges, this work proposes a novel network traffic anomaly detection model for resource-constrained networks based on network flows and WiSARD weightless neural network. The proposed model comprises a novel network flow encoding scheme, WiSARD hyperparameter calibration process, and evaluation strategies tailored for the IDS context. The model is evaluated using real-world traffic data from Rede Rio, as well as public datasets widely adopted in IDS benchmarking studies. Performance comparisons are conducted against established baseline models. Experimental results demonstrate that the proposed approach is capable of effectively detecting network threats, highlighting its potential as a viable solution for flow-based intrusion detection, especially in scenarios involving unknown or emerging attacks.

This work is divided as the following....